\section*{3. Adversarial Alien Obfuscation}

\subsection*{Methodology}
I evaluated zero-shot robustness by replacing alien operator names with different words following four obfuscation schemes:
\begin{itemize}
    \item \textbf{Semantic}: Idea-related English words (e.g., \texttt{combine}).
    \item \textbf{Misleading}: Unrelated English words (e.g., \texttt{paint}).
    \item \textbf{Gibberish}: Random strings (e.g., \texttt{qahft}).
    \item \textbf{Unicode}: Non-ASCII symbols (e.g., $\Delta$).
\end{itemize}

\subsection*{Results \& Analysis}
Table \ref{tab:obfuscation_results} shows the results.

\begin{table}[h]
\centering
\begin{tabular}{l c c c c}
\toprule
\textbf{Scheme} & \textbf{Overall Acc.} & \textbf{Level 1} & \textbf{Level 2} & \textbf{Level 3} \\
\midrule
Original & 11.3\% & 30.0\% & 2.0\% & 2.0\% \\
Gibberish & 10.0\% & 28.0\% & 0.0\% & 2.0\% \\
Unicode & 6.0\% & 16.0\% & 0.0\% & 2.0\% \\
Semantic & 0.0\% & 0.0\% & 0.0\% & 0.0\% \\
Misleading & 0.0\% & 0.0\% & 0.0\% & 0.0\% \\
\bottomrule
\end{tabular}
\caption{Model accuracy under various obfuscation schemes.}
\label{tab:obfuscation_results}
\end{table}

Both \textbf{Semantic} and \textbf{Misleading} schemes completely failed (succeeded?) (0.0\% accuracy) due to a \textbf{formatting collapse}. While meaningless symbols (Original, Gibberish, Unicode) forced the model to treat the prompt as an abstract logic puzzle and output exactly \texttt{Final Answer: <number>}, real English words triggered pre-trained conversational priors. The model output natural sentences (e.g., \texttt{The final answer is 23.}), causing widespread extraction script failures. Thus, English-word obfuscation actively degrades performance by breaking instruction-following behavior.
